% =====================================================================
%  dsfigures.tex
%  Drawing support for the Data Science notes.
%
%  Loaded from main.tex after ../.shared/utpnotes. Two jobs:
%    1  pgfplots libraries and defaults the generated charts need
%    2  the TikZ vocabulary the hand-drawn structure diagrams use
%
%  The charts themselves live in figures/ and are produced by
%  tools/make-figures.py from the same numbers the chapters' R code
%  uses. Regenerate them with:
%
%      python3 tools/make-figures.py
% =====================================================================

% ---------------------------------------------------------------------
%  1.  PGFPLOTS
% ---------------------------------------------------------------------
%  statistics gives boxplot prepared; groupplots gives the multi-panel
%  figures used for Anscombe's quartet and the correlation gallery.

\usepgfplotslibrary{statistics}
\usepgfplotslibrary{groupplots}

% A course-local violet, for the third group in the box plot. The shared
% palette stops at five accents and the box plot needs a sixth.
\definecolor{DSPurple}{HTML}{6B4FA0}
\colorlet{UTPPurple}{DSPurple}
\colorlet{UTPGoldAlias}{UTPGold}

% ---------------------------------------------------------------------
%  2.  STRUCTURE DIAGRAM STYLES
% ---------------------------------------------------------------------
%  R objects drawn as boxes of cells. An atomic vector is a row of
%  `cell`; a list is a row of `slot`, each holding something of a
%  different shape; a data frame is a row of columns.

\tikzset{
  cell/.style={draw=UTPBlue,fill=UTPPaleBlue,minimum width=8mm,minimum height=7mm,
               font=\scriptsize\color{UTPNavy},inner sep=1pt},
  cellalt/.style={cell,fill=UTPPaleOrange,draw=UTPOrange},
  cellghost/.style={cell,fill=UTPPaleGray,draw=UTPMuted!50,
                    font=\scriptsize\itshape\color{UTPMuted}},
  cellout/.style={cell,fill=UTPPaleTeal,draw=UTPTeal},
  cellbad/.style={cell,fill=UTPPaleRed,draw=UTPRed},
  slot/.style={draw=UTPTeal,fill=UTPPaleTeal,minimum width=17mm,
               minimum height=9mm,align=center,
               font=\scriptsize\color{UTPNavy},inner sep=1.5pt},
  dscol/.style={draw=UTPNavy!45,fill=white,minimum width=17mm,
                minimum height=7mm,font=\scriptsize\color{UTPInk},inner sep=1pt},
  dshead/.style={dscol,fill=UTPPaleBlue,font=\scriptsize\bfseries\color{UTPNavy}},
  tag/.style={font=\scriptsize\color{UTPMuted},inner sep=1.5pt},
  tagb/.style={font=\scriptsize\bfseries\color{UTPNavy},inner sep=1.5pt},
  rarrow/.style={-{Stealth[length=2mm]},semithick,draw=UTPTeal},
  rarrowmuted/.style={-{Stealth[length=2mm]},semithick,draw=UTPMuted},
  % flowchart pieces for the control-flow chapter
  fcbox/.style={draw=UTPBlue,fill=UTPPaleBlue,rounded corners=1mm,
                minimum width=22mm,minimum height=7mm,
                font=\scriptsize\color{UTPNavy},align=center},
  fcdiamond/.style={diamond,aspect=2.2,draw=UTPOrange,fill=UTPPaleOrange,
                    inner sep=1pt,font=\scriptsize\color{UTPNavy},align=center},
  fcstart/.style={rounded rectangle,draw=UTPMuted,fill=UTPPaleGray,
                  font=\scriptsize\color{UTPInk},inner sep=2pt},
  fcedge/.style={-{Stealth[length=2mm]},semithick,draw=UTPMuted},
  fclabel/.style={font=\tiny\color{UTPMuted},inner sep=1.5pt},
}

% A short explanatory line inside a figure frame, under the drawing.
\newcommand{\fignote}[1]{\par\smallskip{\scriptsize\color{UTPMuted}#1\par}}

% Inline R code, so chapters can stop writing \texttt{} by hand.
\providecommand{\rcode}[1]{\texttt{#1}}
